% ═════════════════════════════════════════════════════════════════════════════
%  第一章  引言
% ═════════════════════════════════════════════════════════════════════════════
\section{引言}
\label{sec:intro}

% ─────────────────────────────────────────────────────────────────────────────
\subsection{研究背景与动机}
\label{sec:background}

航空发动机是现代航空工业的核心动力装置，其装配过程涉及数千个零部件、
数百道工序以及严格的公差配合约束，是典型的知识密集型、高精度制造场景。
支撑装配作业的知识分散在三类异构数据源中：
\textbf{维修手册}以 PDF 文档形式承载工序逻辑、工具要求和技术规范等过程性知识；
\textbf{BOM（物料清单）表}以结构化电子表格形式提供零件编号、名称和数量等身份信息；
\textbf{三维CAD模型}以几何数据形式编码零件装配层级关系、配合约束和几何公差等空间结构知识。
三类数据源各有其信息边界：手册不含装配结构信息，BOM 不含装配步骤，CAD 不含操作文字描述——
单一数据源均无法构成完整的装配知识底座。
知识图谱（Knowledge Graph）作为结构化知识表示的核心工具，
以三元组形式编码实体间的语义关系，已被证明在工业领域知识管理中具有重要价值
\cite{hogan2021kg,ji2022kgsurvey}。
如何融合三源异构知识，并以自然语言交互方式辅助装配工人完成精密操作，
是提升航空制造智能化水平的重要课题。

近年来，大语言模型（Large Language Model，LLM）在自然语言理解与生成任务上取得了重大突破，
但其固有的三类局限制约了在专业领域的可靠应用：
\textbf{参数知识幻觉}（生成听似合理的错误事实）、
\textbf{训练知识截止}（无法获取最新领域数据），
以及\textbf{专业推理能力不足}（在涉及精确数值和多跳逻辑的专业领域问题上表现不稳定）
\cite{zhao2023llmsurvey}。
检索增强生成（Retrieval-Augmented Generation，RAG）技术为上述问题提供了系统性解决路径
\cite{lewis2020rag}：
通过在推理时动态检索外部知识库，
将相关文本片段注入LLM上下文窗口，
使模型的生成过程以外部事实为锚，
而非单纯依赖参数记忆。
RAG技术自提出以来经历了从Naive RAG到Advanced RAG再到GraphRAG的快速演进
\cite{gao2023ragsurvey}，
在医疗、法律、金融等专业领域已取得显著进展。

然而，现有RAG方法多针对单一文本知识源设计，
难以满足航空发动机装配领域的以下三类特殊需求：

\begin{enumerate}
  \item \textbf{精确标识符匹配}：
        零件编号（如 PN: 1002-AX）、公差数值（如"0.05\textasciitilde0.12\,mm"）、
        扭矩规格等精确标识符对检索准确性至关重要，
        而单纯的密集向量检索因词汇平滑而弱化此类精确匹配能力；
  \item \textbf{结构化推理需求}：
        零件组成关系、装配工序依赖、配合约束等知识天然具有图结构，
        需要多跳关系推理，远超向量相似度检索的能力边界；
  \item \textbf{跨源一致性约束}：
        手册、BOM 与 CAD 对同一零件或工序常采用不同命名方式，
        若缺乏统一实体注册与跨源对齐，系统容易在装配层级、工具要求和技术规范之间产生断裂；
        此外，检索噪声（上下文中混入的不相关段落）会导致LLM生成质量显著下降
        \cite{liu2023lostinmiddle}，
        对高精度装配场景危害尤甚。
\end{enumerate}

现有工作已开始关注RAG方法在工业故障诊断与装配引导场景中的应用。
在航空故障诊断领域，Tang等人\cite{tang2023aircraft}和Meng等人\cite{meng2024aircraft}
分别系统综述并实现了面向航空故障诊断的知识图谱平台；
Wu等人\cite{wu2024engine}构建了专门面向航空发动机润滑系统的故障KG，
证明了KG在航空发动机专业场景下的工程价值。
Liu等人\cite{liu2024jointkg}在航空装配功能测试场景下，
将知识图谱通过前缀微调嵌入大语言模型，实现了98.5\%的故障诊断准确率，
证明了专业领域KG与LLM协同推理的有效性。
Liu等人\cite{liu2023avkg}构建了含5类实体、3类关系共1308组三元组的
航空装配知识图谱，并与Neo4j图数据库集成，为装配知识管理奠定了基础。
Shi等人\cite{shi2022arkg}提出基于"功能-结构-工序-装配资源"多维本体的
装配资源知识图谱（ARKG），将CAD模型数据经本体映射后融入知识图谱，
支持工程决策的知识共享与复用。
Shen等人\cite{shen2025hybridrag}提出了HybridRAG框架，
将知识图谱（KG）与向量检索融合，
并在飞机故障排查场景中验证了混合检索相较于单一稠密检索路径的优越性，
F1提升约4\%、幻觉率降低约7\%。
然而，上述工作多以手册文本或 CAD 属性的单一侧面为主，
尚未将 BOM 表的权威零件身份信息、CAD/IPC 装配结构信息与维修手册工序知识系统融入同一图谱。
这使其在回答"某零件属于哪个上级装配体"、"某工序需要哪些工具和规范"等问题时，
容易出现实体缺失、关系链断裂和跨章节召回不稳定等现象。
现有方法在装配这一三源知识协同的精细场景下仍存在明显局限。

% ─────────────────────────────────────────────────────────────────────────────
\subsection{现有方法局限性}
\label{sec:limitation}

Fan等人\cite{fan2024ragsurvey}的综述指出，RAG方法的核心挑战在于检索质量、
增强策略和生成可靠性三个维度；Gao等人\cite{gao2023ragsurvey}进一步将RAG发展
划分为Naive RAG、Advanced RAG和GraphRAG三个演进阶段，指出GraphRAG在结构化推理
方面具有显著优势但知识来源仍局限于文本。
结合航空发动机装配的特殊需求，本文进一步识别出以下三方面不足：

\textbf{（1）缺乏稀疏-稠密互补机制。}
大多数方案选择单一的稠密向量检索路径，
在处理含精确零件编号和数字标识符的领域特定查询时召回率偏低。
BM25等稀疏检索方法虽对词汇精确匹配有天然优势，
但单独使用时语义泛化能力不足，
两类方法的互补融合尚未在装配领域得到充分研究。

\textbf{（2）三源异构知识融合不足。}
现有工业知识图谱工作多从维修手册文本单一来源提取三元组：
Liu等人\cite{liu2023avkg}的航空装配KG仅含5类实体和3类关系，
Shi等人\cite{shi2022arkg}的ARKG虽引入CAD模型数据，但未实现BOM、CAD、手册三者的
系统融合与实体对齐。三源数据的互补边界清晰
（手册提供过程知识，BOM提供身份信息，CAD提供几何拓扑），
但现有方法缺乏对三者进行统一本体建模和跨源实体对齐的完整框架。

\textbf{（3）多源异质结果融合方法简单。}
多路检索的结果评分处于不同尺度空间，
现有方法多采用线性加权融合（$\text{IR} = \alpha \cdot \text{score}_{\text{dense}} + \beta \cdot \text{score}_{\text{sparse}}$）
\cite{shen2025hybridrag}，
易受评分分布差异影响，
缺乏对文本块与图谱子图等异质内容统一排序的无参融合机制。

% ─────────────────────────────────────────────────────────────────────────────
\subsection{本文主要贡献}
\label{sec:contribution}

针对上述不足，本文提出面向航空发动机装配的
\textbf{多源混合检索增强生成}框架
（Multi-source Hybrid RAG，MHRAG），
主要贡献如下：

\begin{enumerate}
  \item \textbf{三源知识图谱构建方法}。
        设计了以BOM解析、CAD结构化解析、LLM手册三元组抽取为三轨的
        多智能体协同知识抽取流水线，三轨结果通过实体对齐后
        汇入统一的Neo4j图数据库。
        本体模式包含7类实体（含 \textit{GeoConstraint} 几何约束节点）
        和9类关系，清晰界定了三源数据的语义边界：
        BOM提供零件身份注册表，CAD提供装配层级与几何约束，
        手册提供工序过程与技术规范。
        知识融合阶段引入验证智能体进行一致性校验，
        有效解决同名异物和信息冲突问题。

  \item \textbf{三路并行混合检索架构}。
        在密集向量检索（BGE-M3）和BM25稀疏检索之外，
        引入基于BGE-M3语义近邻搜索（ANN）的知识图谱子图召回路径，
        使系统同时具备语义召回、精确字符串匹配和结构化关系补全能力。
        三路结果通过互惠排名融合（RRF）统一排序，
        无需对各路评分进行归一化处理，对单路检索失败具有天然鲁棒性；
        相比传统的关键词字符串匹配（CONTAINS）KG检索方法，
        语义ANN方案更适合处理中英混合查询与专业术语别名。

  \item \textbf{面向装配引导的生成与交互设计}。
        引入多查询扩展策略缓解词汇不匹配问题；
        采用基于BERT的交叉编码器精排进一步提升上下文质量；
        设计结构化系统提示使LLM以"逐步确认"模式输出操作指令，
        符合航空装配的安全性要求；
        答案通过SSE流式返回并附带来源引用，
        支持前端展示可追溯的手册片段与图谱关系。

  \item \textbf{系统实现与验证}。
        基于FastAPI + Neo4j + Qdrant实现了完整的后端服务，
        前端以React实现交互式问答界面，
        在 PT6A 整机维修手册（10 个 ATA 章节、约 1.78\,MB Markdown 文本）上构建出
        含 7{,}333 节点 / 8{,}752 边的知识图谱。
        在含 60 题精评（72-30 压气机子系统）与 20 题跨章节粗评的双层 QA 评测集上，
        本文方法（MHRAG）在 Tier-1 精评的 LLM-Judge Relevance、Completeness 与
        Correctness 维度取得 3.50、3.02 与 2.35 的表现；
        在 Tier-2 跨章节粗评中 Completeness 与 Correctness 维度领先所有单路基线
        （相对 BM25-only 分别提升 8.9\% 与 1.9\%）；
        消融实验进一步分析各路径的独立贡献：
        KG 路径主要提升答案完整性与事实一致性，
        BM25 路径则保证零件编号、扭矩和公差等精确表达的可召回性。
        工程性能维度，本文方法在 7.76\,s 端到端延迟和 1.60\,GB 本地显存占用下
        取得 35.00\% 的答非所问率，
        契合消费级 GPU 的工业部署约束。
\end{enumerate}

% ─────────────────────────────────────────────────────────────────────────────
\subsection{论文结构}
\label{sec:structure}

本文其余部分组织如下：
第\ref{sec:related}节回顾RAG、知识图谱构建与混合检索的相关工作；
第\ref{sec:method}节详细介绍本文提出的MHRAG框架各组件；
第\ref{sec:experiment}节描述实验配置、数据集构建与评价指标，
并汇报与基线方法的对比结果及消融分析；
第\ref{sec:discussion}节讨论各检索路径的贡献、知识图谱质量影响以及方法局限性；
第\ref{sec:conclusion}节总结全文并展望未来工作方向。
